\documentclass[11pt]{article}
\usepackage[T1]{fontenc}
\usepackage{lmodern}
\usepackage{microtype}
\usepackage{amsmath,amssymb,amsthm,mathtools}
\usepackage{graphicx}
\usepackage[margin=1in]{geometry}
\usepackage[hidelinks]{hyperref}
\usepackage{xcolor}

\newtheorem{theorem}{Theorem}
\newtheorem{lemma}{Lemma}
\newtheorem{remark}{Remark}
\newcommand{\A}{\mathcal A}
\newcommand{\B}{\mathcal B}
\newcommand{\one}{\mathbf 1}
\newcommand{\Ncap}[1]{\lVert #1\rVert_{1\cap\infty}}
\newcommand{\Nsum}[1]{\lVert #1\rVert_{1+\infty}}

\title{Duality closes the intermediate-trace gap for\\
positive isometries of $L^1\cap L^\infty$}
\author{Candidate full solution packet for arXiv:2007.02324}
\date{27 August 2026}

\begin{document}
\maketitle

\noindent\textbf{Status.}
Candidate full solution, likely valid, pending expert review.  The argument
answers the open problem in Remark~5.6 of de Jager--Conradie
\cite{deJagerConradieExtreme} and in fact proves a slightly stronger theorem.

\section{Source question}

The source is Pierre de Jager and Jurie Conradie, \emph{Extreme point methods
in the study of isometries on certain non-commutative spaces},
arXiv:2007.02324, Remark~5.6 on printed/PDF page~21
\cite{deJagerConradieExtreme}.  The complete source passage is reproduced
below directly from the paper.

\begin{center}
\includegraphics[width=0.97\textwidth]{figures/open_problem_crop.png}
\end{center}

\noindent\textbf{Transcription of the question.}
Describe the positive surjective isometries
\[
 U:L^1\cap L^\infty(\A)\longrightarrow L^1\cap L^\infty(\B)
\]
if $1<\tau(\one_\A)<2$ or $1<\nu(\one_\B)<2$.  Here $(\A,\tau)$ and
$(\B,\nu)$ are non-atomic semifinite von Neumann algebras, as throughout the
source section.

For a semifinite traced algebra $(\mathcal M,\mu)$, write
\[
 E_\mathcal M=L^1(\mathcal M,\mu)\cap\mathcal M,
 \qquad
 \|x\|_{1\cap\infty}=\max\{\|x\|_1,\|x\|_\infty\},
\]
and
\[
 F_\mathcal M=L^1(\mathcal M,\mu)+\mathcal M,
 \qquad
 \|a\|_{1+\infty}=
 \inf_{a=a_1+a_\infty}
   \bigl(\|a_1\|_1+\|a_\infty\|_\infty\bigr).
\]

\section{Proof intuition}

The obstruction noted by the source is specific to its direct extreme-point
argument: when the total trace is below two, there are not two orthogonal
trace-one projections.  The key idea is to change sides of the duality rather
than repair that argument.

If $1<\tau(\one_\A)<2$, then $E_\A=\A$ and $\one_\A$ is an order unit.  Since
a positive onto isometry is an order isomorphism, its image is an order unit
in $E_\B$.  An infinite-trace intersection space has no order unit; hence
\(\nu(\one_\B)<\infty\) and $E_\B=\B$.  Extreme points then rule out
\(\nu(\one_\B)\leq1\).

Both algebras are therefore finite with total trace greater than one.  The
normal preadjoint of $U$ is a positive surjective isometry for the polar
\(L^1+L^\infty\) norms.  The source paper has already classified precisely
those sum-space isometries when both total traces exceed one.  It gives a
trace-preserving Jordan isomorphism, and the trace pairing shows that $U$ is
its inverse.

\section{Result}

\begin{theorem}[Intermediate-trace classification]\label{thm:main}
Let $(\A,\tau)$ and $(\B,\nu)$ be non-atomic semifinite von Neumann algebras,
and let
\[
 U:E_\A\longrightarrow E_\B
\]
be a positive surjective linear isometry.  Suppose that
\(1<\tau(\one_\A)<\infty\) or \(1<\nu(\one_\B)<\infty\).  Then both identity
traces are finite and greater than one, they are equal, and $U$ is a
trace-preserving Jordan $*$-isomorphism from $\A$ onto $\B$.

Conversely, every trace-preserving Jordan $*$-isomorphism from $\A$ onto
$\B$ restricts to a positive surjective isometry $E_\A\to E_\B$.
In particular, this completely answers the source question, whose finite
trace is assumed to lie in $(1,2)$.
\end{theorem}

\begin{proof}
De Jager--Conradie prove that a positive isometry between noncommutative
symmetric spaces reflects positivity; if it is onto, it and its inverse are
order isomorphisms and are normal
\cite[Lemma~3.1]{deJagerConradieSymmetric}.  Thus, if the finite trace in the
hypothesis is $\nu(\one_\B)$, we may replace $U$ by $U^{-1}$.  Assume from now
on that
\[
 1< t:=\tau(\one_\A)<\infty.
\]
Then $E_\A=\A$ as a set, and $\one_\A$ is an order unit.

\medskip
\noindent\emph{Step 1: the target trace is finite.}
Set $y=U(\one_\A)\geq0$.  For every $z\in(E_\B)_{\rm sa}$, positivity of
$U$ and $U^{-1}$ gives
\[
 -\|U^{-1}z\|_\infty y\leq z\leq
 \|U^{-1}z\|_\infty y.
\]
Hence $y$ is an order unit of $E_\B$.

We claim that $E_\B$ has no order unit when $\nu(\one_\B)=\infty$.  Indeed,
let $0\leq y\in E_\B$.  For each integer $k\geq1$, Markov's inequality gives
\[
 \nu\!\left(\mathbf1_{(1/k,\infty)}(y)\right)\leq k\nu(y)<\infty,
\]
so $e_k=\mathbf1_{[0,1/k]}(y)$ has infinite trace.  By semifiniteness and
non-atomicity, recursively choose nonzero, mutually orthogonal projections
$q_k\leq e_k$ with $\nu(q_k)\leq2^{-k}$.  The strong sum
$q=\sum_{k\geq1}q_k$ is a projection of finite trace, so $q\in E_\B$.
If $q\leq cy$ for some $c>0$, compression by $q_k$ yields
\[
 q_k\leq c q_k yq_k\leq \frac{c}{k}q_k,
\]
which is impossible for $k>c$.  Thus $y$ is not an order unit.  This
contradiction proves $s:=\nu(\one_\B)<\infty$, and hence $E_\B=\B$.

\medskip
\noindent\emph{Step 2: the target total is greater than one.}
The source paper characterizes the extreme points of the intersection-space
unit ball: they are the partial isometries $v$ satisfying
\[
 \nu(|v|)=\min\{1,\nu(\one_\B)\}
\]
(and analogously for $\A$)
\cite[Proposition~5.1]{deJagerConradieExtreme}.  Since $\A$ is non-atomic and
$t>1$, it has two distinct projections $p_1,p_2$ of trace one.  (Starting from
one such projection, move equal small-trace pieces between it and its nonzero
complement.)  Both $p_i$ are extreme points.

A surjective linear isometry preserves extreme points, so each $U(p_i)$ is a
positive extreme point.  A positive partial isometry is a projection.  If
$s\leq1$, the displayed characterization forces
\(\nu(U(p_i))=s=\nu(\one_\B)\), whence faithfulness gives
$U(p_i)=\one_\B$.  This contradicts injectivity because $p_1\ne p_2$.
Therefore $s>1$.

\medskip
\noindent\emph{Step 3: pass to the normal dual.}
On each finite algebra the intersection norm is equivalent to the operator
norm:
\[
 \|x\|_\infty\leq\Ncap{x}\leq
 \max\{1,\mu(\one)\}\|x\|_\infty.
\]
Consequently the positive normal maps $U:\A\to\B$ and $U^{-1}$ have inverse
normal preadjoints.  Write
\begin{equation}\label{eq:pairing}
 V:L^1(\B,\nu)\longrightarrow L^1(\A,\tau),
 \qquad \tau(V(b)x)=\nu(bU(x)).
\end{equation}
For normal functionals, intersection/sum polarity gives
\begin{equation}\label{eq:polar}
 \|a\|_{1+\infty}
 =\sup\bigl\{|\mu(ax)|:\ x\in\mathcal M,
                   \ \|x\|_{1\cap\infty}\leq1\bigr\}.
\end{equation}
This is the standard noncommutative K\"othe-duality formula
\cite{DoddsDoddsDePagter}.  Equivalently, it is the polar formula for the
maximum norm on an intersection:
the polar unit ball is the closed convex hull of the two polar unit balls,
equivalently the infimal-convolution sum norm.  It remains valid here after
restricting to normal functionals, which are represented by $L^1$.

Using \eqref{eq:pairing}, \eqref{eq:polar}, and
$U(B_{E_\A})=B_{E_\B}$, we obtain
\[
 \|V(b)\|_{1+\infty}
 =\sup_{\Ncap{x}\leq1}|\nu(bU(x))|
 =\sup_{\Ncap{z}\leq1}|\nu(bz)|
 =\|b\|_{1+\infty}.
\]
The preadjoint of $U^{-1}$ is $V^{-1}$, and positivity of $U$ and $U^{-1}$
implies positivity of both preadjoints.  Thus
\[
 V:F_\B\longrightarrow F_\A
\]
is a positive surjective isometry.  Notice that finite trace makes
$F_\A=L^1(\A,\tau)$ and $F_\B=L^1(\B,\nu)$ as sets, though not with their
ordinary $L^1$ norms.

\medskip
\noindent\emph{Step 4: apply the sum-space classification and return.}
Both total traces are now greater than one, so Theorem~4.2 and its positive
corollary in the source paper apply to $V$
\cite[Theorem~4.2]{deJagerConradieExtreme}.  They show that $V$ is the
restriction of a trace-preserving Jordan $*$-isomorphism
\[
 J:\B\longrightarrow\A.
\]
We finally identify the original map.  Jordan isomorphisms preserve
$a\circ b=(ab+ba)/2$, and a trace satisfies
$\mu(ab)=\mu(a\circ b)$.  Therefore, for $b\in\B$ and $x\in\A$,
\begin{equation}\label{eq:trace-pairing}
 \tau(J(b)x)
 =\tau(J(b)\circ x)
 =\tau\!\left(J\bigl(b\circ J^{-1}(x)\bigr)\right)
 =\nu\bigl(b\circ J^{-1}(x)\bigr)
 =\nu(bJ^{-1}(x)).
\end{equation}
Equations \eqref{eq:pairing}, \eqref{eq:trace-pairing}, and $V=J$ give
\(\nu(bU(x))=\nu(bJ^{-1}(x))\) for all $b\in\B$.  Since bounded normal
functionals separate points of the finite algebra $\B$, $U(x)=J^{-1}(x)$ for
all $x\in\A$.  Hence $U$ is a trace-preserving Jordan $*$-isomorphism.
In particular it is unital, so
\[
 \tau(\one_\A)=\nu(\one_\B).
\]

The converse is exactly the norm calculation in the source theorem: a
trace-preserving Jordan $*$-isomorphism preserves both the $L^1$ norm and the
operator norm, hence their maximum.  This completes the proof.
\end{proof}

\section{Verification, limitations, and novelty status}

The proof is non-computational.  Its two external mathematical inputs were
checked in the locally stored papers: the order-isomorphism/normality result is
Lemma~3.1 of \cite{deJagerConradieSymmetric}, and the sum-space classification
is Theorem~4.2 with its positive corollary in
\cite{deJagerConradieExtreme}.  The source paper's Proposition~5.1 supplies
the extreme-point characterization used only to exclude a target total at
most one.

The theorem addresses positive surjective linear isometries under the source
paper's standing non-atomicity and semifiniteness assumptions.  It does not
classify non-positive isometries, nor does it remove non-atomicity.

A bounded literature search through 27 August 2026 used arXiv:2007.02324, the
exact title and authors, the quoted open-problem phrase, combinations of
``positive surjective isometry'', ``$L^1\cap L^\infty$'', and ``trace between
1 and 2'', and a forward-citation check.  It found the arXiv/journal versions
of the source and one unrelated citing paper, but no later claimed answer to
this question.  This supports, but does not guarantee, novelty.

\medskip
\noindent\textbf{Human-review recommendation.}
Review as a candidate full result.  The main checkpoint is formula
\eqref{eq:polar} on the normal predual and its use to make the preadjoint an
onto sum-space isometry.  The order-unit and trace-pairing arguments should be
checked next.

\begin{thebibliography}{9}

\bibitem{deJagerConradieExtreme}
P. de Jager and J. Conradie,
\emph{Extreme point methods in the study of isometries on certain
non-commutative spaces}, arXiv:2007.02324; Glasgow Math. J.,
doi:10.1017/S0017089521000227.

\bibitem{deJagerConradieSymmetric}
P. de Jager and J. J. Conradie,
\emph{Isometries between non-commutative symmetric spaces associated with
semi-finite von Neumann algebras}, Positivity \textbf{24} (2020), 815--835;
arXiv:1907.06452.

\bibitem{DoddsDoddsDePagter}
P. G. Dodds, T. K.-Y. Dodds, and B. de Pagter,
\emph{Non-commutative K\"othe duality}, Trans. Amer. Math. Soc.
\textbf{339} (1993), 717--750.

\end{thebibliography}

\end{document}
